\scanpage{39}
\uline{Th. 1} \inserted{Soit $n\geq1$.} Le foncteur $E_n^\circ$ passe à la catégorie isotopique $f_n^\circ$, et induit un foncteur \uuline{pl. fidèle}
\[ f_n^\circ\longrightarrow(G_n\text{-ens}) \]
(dont l'image essentielle est formée des $G_n$-ens. tels que les $\sigma_i$ $(0\leq i\leq n)$ \inserted{y} opèrent \deleted{sans pts fixes} \inserted{librement}.

Considérons $X_n\in\mathrm{Ob}\,\mathcal F_n$. On lui associe une cascade d'ens. à opérateurs
\diagram{$\begin{array}{lll} R_n(X_n)\overset{\mathrm{dfn}}{=}E_n(X_n)=\unsure{1|I}&G_{n-1}\text{-ens}&\sigma_0\ldots\sigma_{n-1}\\ R_{n-1}(X_n)=E_{n-1}(X_{n-1})&G_{n-2}\text{-ens}&\sigma_0\ldots\sigma_{n-2}\\ \ldots&&\\ R_1(X_n)=E_1(X_1)&G_0\text{-ens}&\sigma_0\\ R_0(X_n)=E_0(X_0)\overset{\mathrm{dfn}}{=}X_0&G_{-1}\ \text{ens}&\text{i.e. ensemble}\end{array}$}
De plus on a \deleted{\unsure{de}} \deleted{\unsure{des appl.}} \inserted{une} cascade d'applications
\[ R_n(X_n)\to R_{n-1}(X_n)\to\ldots\to R_0(X_n) \]
\deleted{\unsure{définie|définies}}
\diagram{\begin{tikzcd}[ampersand replacement=\&,row sep=.15em]R_i(X_n)\arrow[r]\&R_{i-1}(X_n)\&\text{hom de }G_{i-2}\text{-ens}\\\rotatebox{90}{$=$}\&\rotatebox{90}{$=$}\&\text{(commute à}\\E_i(X_i)\&E_{i-1}(X_{i-1})\&\sigma_0\ldots\sigma_{i-2})\\\rotatebox{90}{$=$}\&{}\&{}\\E_{i-1}^\circ(\partial\widetilde X_i)=E_{i-1}(\partial\widetilde X_i)\&{}\&{}\end{tikzcd}}
se déduit de $\partial\widetilde X_i\unsure{\text{---}|\to}X_{i-1}$ \unsure{en|par} appliquant le foncteur $E_{i-1}$.

\uline{Th 2} \inserted{Soit $n\geq0$.} Soit $(\mathrm{Casc}_n)$ la catégorie des ``cascades d'ens à op'' sous les grpes $G_{\unsure{n-1|n+1}},G_n,\ldots,G_0$, $G_{-1}=\{1\}$, on a donc un foncteur
% archive sheet number: 38
